% !TEX root = ../thesis-main.tex

\chapter{Introduction}
\label{C:introduction}

\begin{epigraph}%[Wallace Stevens]
  I was of three minds, \\
  Like a tree \\
  in which there are three blackbirds.\footnotemark
\end{epigraph}
\footnotetext{The chapter epigraphs are stanzas II, XII, IX, III, and XIII respectively, from the poem ``Thirteen Ways of Looking at a Blackbird'' by Wallace Stevens (1879-1955) \cite{Stevens1923}.}

This thesis consists of three parts: the front matter, which you have just passed, the \textit{kappa}, which you are reading, and the included papers, which you will get to shortly. This thesis is also the outcome of approximately three research questions: one about computing relative homological algebra for poset-indexed functors, another about inducing barcode matchings from morphisms of persistence modules, and a third about generalizing Morse theory to distance functions between algebraic varieties. All of these questions, their answers, and the resulting further questions will be presented in this \textit{kappa}. Most importantly, this thesis revolves around three mathematical trios: filtered topological, algebraic, and combinatorial spaces; algebraic-topological, -geometric, and -\nobreak persistent invariants; and generic, algorithmic, and numeric computations.

There is no one-to-one correspondence between the research questions and the mathematical concepts; each project involves each concept to some extent. However, the results of this thesis can be split into those studying persistence modules and those studying Morse theory.

In the words of Gunnar Carlsson, a foundational researcher in \index{topological data analysis (TDA)}\textbf{topological data analysis} (\textbf{TDA}), ``data has shape and the shape matters'' \cite{Carlsson2015}. TDA is a relatively new field, established around the turn of the $\text{21}^{\text{st}}$ century, that applies invariants from topology and algebraic topology to understand the topological and geometric structure of data. One of the central objects in TDA, which quantifies the shape of, and encodes information about, data is the \textbf{persistence module}, which is the image of a homology functor applied to a filtered topological space. In \cref{C:persistent-algebra}, we present \textbf{persistent algebra}, the study of persistence modules, which will be relevant for Papers \ref{PC:CGRST2024}, \ref{PC:AGRS2025}, and \ref{PC:GRRS2026}.

While the first two-thirds of this thesis revolve around persistence modules, the last third studies filtered spaces form a different angle, namely that of \index{Morse theory}\textbf{Morse theory}. Morse theory was first developed by the eponymous mathematician under the name \textit{calculus of variations in the large} \cite{Morse1934}. In \cref{C:metric-morse} we give a modern presentation of the original theory for smooth manifolds, as well as its extensions to continuous selections, which are the basis for Papers \ref{PC:GLRS2025} and \ref{PC:Ren2026}.

In \cref{C:contributions} we summarize the results of each included paper, and in \cref{C:outlook} we present some ongoing work and future directions of research.